\documentclass[11pt,a4paper]{article}
\usepackage[T1]{fontenc}
\usepackage[utf8]{inputenc}
\usepackage{lmodern,microtype}
\usepackage[margin=2.2cm]{geometry}
\usepackage{booktabs,tabularx,array}
\usepackage{xcolor,amssymb,enumitem}
\usepackage{forest,adjustbox}
\usepackage{caption}
\usepackage[hidelinks]{hyperref}
\setlist{itemsep=1pt,topsep=3pt}
\captionsetup{font=small}

\newcommand{\tg}[1]{\textsubscript{\sffamily #1}}
\makeatletter\newcommand{\tinput}[1]{\@@input #1 }\makeatother % \input usable inside tables

\begin{document}

\begin{center}
  {\Large\bfseries NLP Today -- Session 1: PoS Tagging and Constituency Parsing}\\[8pt]
  % ======================= EDIT HERE: group members =======================
  \textcolor{red}{[First name LAST NAME] -- [e-mail]}\\
  \textcolor{red}{[First name LAST NAME] -- [e-mail]}\\
  \textcolor{red}{[First name LAST NAME] -- [e-mail]}
  % =========================================================================
\end{center}

\noindent\textit{Use of AI tools.} This report was prepared with Claude (Anthropic), an AI assistant, which we used for the code, the experiments and the writing. \textcolor{red}{[Optional: what the group checked or changed.]}

\medskip
\noindent We analyse the six sentences of the dataset: four general-domain sentences and two scientific sentences (biology and astronomy). Part~1 gives our manual analysis, Part~2 the automatic analysis (PoS tagging with spaCy and sciSpaCy, constituency parsing with the CYK algorithm and with the Berkeley Neural Parser), and Part~3 our observations.

%==========================================================================
\section{Manual analysis}

\subsection{Part-of-speech tags}
We split the sentences on spaces and punctuation, but kept \emph{Dec.}, \emph{19.33}, \emph{0.044}, \emph{type-Ia} and \emph{J10354824+3900279} as single tokens. We used the ten allowed Universal Dependencies tags. Three words (\emph{they}, \emph{which}, \emph{this}) are pronouns and none of the allowed tags fits them, so we tagged them PRON (marked with $\dagger$).

{\small\raggedright\begin{description}[leftmargin=2em, style=sameline, itemsep=1pt]
\item[S1] The\tg{DET} cat\tg{NOUN} sat\tg{VERB} on\tg{ADP} the\tg{DET} couch\tg{NOUN} .\tg{PUNCT}
\item[S2] Time\tg{NOUN} flies\tg{VERB} like\tg{ADP} an\tg{DET} arrow\tg{NOUN} .\tg{PUNCT}
\item[S3] The\tg{DET} spy\tg{NOUN} saw\tg{VERB} the\tg{DET} cop\tg{NOUN} with\tg{ADP} the\tg{DET} telescope\tg{NOUN} .\tg{PUNCT}
\item[S4] The\tg{DET} spy\tg{NOUN} saw\tg{VERB} the\tg{DET} cop\tg{NOUN} with\tg{ADP} the\tg{DET} revolver\tg{NOUN} .\tg{PUNCT}
\item[S5] Arabidopsis\tg{PROPN} thaliana\tg{PROPN} seedlings\tg{NOUN} exhibit\tg{VERB} longer\tg{ADJ} hypocotyls\tg{NOUN} when\tg{SCONJ} they\tg{PRON$^\dagger$} are\tg{AUX} grown\tg{VERB} under\tg{ADP} high\tg{ADJ} ambient\tg{ADJ} temperature\tg{NOUN} ,\tg{PUNCT} which\tg{PRON$^\dagger$} is\tg{AUX} defined\tg{VERB} as\tg{ADP} thermomorphogenesis\tg{NOUN} .\tg{PUNCT}
\item[S6] A\tg{DET} spectrogram\tg{NOUN} of\tg{ADP} PSN\tg{PROPN} J10354824+3900279\tg{PROPN} obtained\tg{VERB} on\tg{ADP} Dec.\tg{PROPN} 19.33\tg{NUM} UT\tg{PROPN} suggests\tg{VERB} that\tg{SCONJ} this\tg{PRON$^\dagger$} is\tg{AUX} a\tg{DET} type-Ia\tg{NOUN} at\tg{ADP} redshift\tg{NOUN} z\tg{NOUN} 0.044\tg{NUM} .\tg{PUNCT}
\end{description}
\par}

\noindent Some choices are debatable: we tagged the species name \emph{Arabidopsis thaliana} and \emph{UT} (Universal Time) as PROPN, and the variable \emph{z} as NOUN.

\subsection{Ambiguities and expected trees}
Out of context, several words can have another PoS than the one we chose: \emph{couch} (VERB), \emph{Time} (VERB), \emph{flies} (NOUN), \emph{like} (VERB), \emph{spy} (VERB), \emph{saw} (NOUN), \emph{cop} (VERB), \emph{telescope} (VERB), \emph{Arabidopsis} (NOUN), \emph{thaliana} (NOUN), \emph{exhibit} (NOUN), \emph{grown} (ADJ), \emph{high} (NOUN), \emph{defined} (ADJ), \emph{as} (SCONJ), \emph{obtained} (ADJ), \emph{UT} (NOUN), \emph{that} (DET, PRON), \emph{this} (DET), \emph{type-Ia} (ADJ)\unskip. The dataset also contains structural ambiguities:
\begin{itemize}
  \item \emph{Time flies like an arrow} can mean that time passes quickly (\emph{flies} is the verb), that ``time flies'' (a kind of fly) like an arrow (\emph{like} is the verb), or it can be an order to time flies (\emph{Time} is the verb). We expect the first reading.
  \item In S3 and S4, the phrase \emph{with the telescope} / \emph{with the revolver} can be attached to the verb \emph{saw} or to the noun \emph{cop}. We expect the verb attachment for S3 (the telescope is used to see) and the noun attachment for S4 (the cop has the revolver), although both sentences have exactly the same PoS tags.
  \item In S5, \emph{which is defined as thermomorphogenesis} refers to the whole situation (seedlings growing longer when it is warm), not only to \emph{high ambient temperature}.
  \item In S6, \emph{obtained on Dec.~19.33 UT} describes the spectrogram, and \emph{at redshift z 0.044} describes the supernova.
\end{itemize}
Figures~\ref{fig:general} and~\ref{fig:scientific} show our expected trees. We only use the labels S, NP, VP and PP (plus the PoS tags). Because the SBAR label is not allowed, subordinate clauses are labelled S. The trees are binary so that they can be produced by the CYK algorithm (Part~2).

\begin{figure}[h]\centering
\begin{minipage}[b]{0.45\linewidth}\centering\begin{adjustbox}{max width=\linewidth}\begin{forest} for tree={s sep=1.2mm, l sep=2pt, l=0, inner sep=1pt, font=\footnotesize}, where n children=0{font=\footnotesize\itshape}{} [{S} [{S} [{NP} [{DET} [{The}, tier=word]] [{NOUN} [{cat}, tier=word]]] [{VP} [{VP} [{sat}, tier=word]] [{PP} [{ADP} [{on}, tier=word]] [{NP} [{DET} [{the}, tier=word]] [{NOUN} [{couch}, tier=word]]]]]] [{PUNCT} [{.}, tier=word]]] \end{forest}\end{adjustbox}\\\small S1\end{minipage}\hfill
\begin{minipage}[b]{0.45\linewidth}\centering\begin{adjustbox}{max width=\linewidth}\begin{forest} for tree={s sep=1.2mm, l sep=2pt, l=0, inner sep=1pt, font=\footnotesize}, where n children=0{font=\footnotesize\itshape}{} [{S} [{S} [{NP} [{Time}, tier=word]] [{VP} [{VP} [{flies}, tier=word]] [{PP} [{ADP} [{like}, tier=word]] [{NP} [{DET} [{an}, tier=word]] [{NOUN} [{arrow}, tier=word]]]]]] [{PUNCT} [{.}, tier=word]]] \end{forest}\end{adjustbox}\\\small S2\end{minipage}\\[8pt]
\begin{minipage}[b]{0.49\linewidth}\centering\begin{adjustbox}{max width=\linewidth}\begin{forest} for tree={s sep=1.2mm, l sep=2pt, l=0, inner sep=1pt, font=\footnotesize}, where n children=0{font=\footnotesize\itshape}{} [{S} [{S} [{NP} [{DET} [{The}, tier=word]] [{NOUN} [{spy}, tier=word]]] [{VP} [{VP} [{VERB} [{saw}, tier=word]] [{NP} [{DET} [{the}, tier=word]] [{NOUN} [{cop}, tier=word]]]] [{PP} [{ADP} [{with}, tier=word]] [{NP} [{DET} [{the}, tier=word]] [{NOUN} [{telescope}, tier=word]]]]]] [{PUNCT} [{.}, tier=word]]] \end{forest}\end{adjustbox}\\\small S3 (PP attached to the verb phrase)\end{minipage}\hfill
\begin{minipage}[b]{0.49\linewidth}\centering\input{gen/target_S4.tex}\\\small S4 (PP attached to the noun phrase)\end{minipage}
\caption{Expected trees for the general-domain sentences.}\label{fig:general}
\end{figure}

\begin{figure}[h]\centering
\begin{adjustbox}{max width=\linewidth}\begin{forest} for tree={s sep=1.2mm, l sep=2pt, l=0, inner sep=1pt, font=\footnotesize}, where n children=0{font=\footnotesize\itshape}{} [{S} [{S} [{S} [{NP} [{PROPN} [{Arabidopsis}, tier=word]] [{NP} [{PROPN} [{thaliana}, tier=word]] [{NP} [{seedlings}, tier=word]]]] [{VP} [{VP} [{VERB} [{exhibit}, tier=word]] [{NP} [{ADJ} [{longer}, tier=word]] [{NP} [{hypocotyls}, tier=word]]]] [{S} [{SCONJ} [{when}, tier=word]] [{S} [{NP} [{they}, tier=word]] [{VP} [{VP} [{AUX} [{are}, tier=word]] [{VERB} [{grown}, tier=word]]] [{PP} [{ADP} [{under}, tier=word]] [{NP} [{ADJ} [{high}, tier=word]] [{NP} [{ADJ} [{ambient}, tier=word]] [{NP} [{temperature}, tier=word]]]]]]]]]] [{S} [{PUNCT} [{,}, tier=word]] [{S} [{NP} [{which}, tier=word]] [{VP} [{VP} [{AUX} [{is}, tier=word]] [{VERB} [{defined}, tier=word]]] [{PP} [{ADP} [{as}, tier=word]] [{NP} [{thermomorphogenesis}, tier=word]]]]]]] [{PUNCT} [{.}, tier=word]]] \end{forest}\end{adjustbox}\\\small S5\\[8pt]
\begin{adjustbox}{max width=\linewidth}\begin{forest} for tree={s sep=1.2mm, l sep=2pt, l=0, inner sep=1pt, font=\footnotesize}, where n children=0{font=\footnotesize\itshape}{} [{S} [{S} [{NP} [{NP} [{NP} [{DET} [{A}, tier=word]] [{NOUN} [{spectrogram}, tier=word]]] [{PP} [{ADP} [{of}, tier=word]] [{NP} [{PROPN} [{PSN}, tier=word]] [{NP} [{J10354824+3900279}, tier=word]]]]] [{VP} [{VP} [{obtained}, tier=word]] [{PP} [{ADP} [{on}, tier=word]] [{NP} [{PROPN} [{Dec.}, tier=word]] [{NP} [{NUM} [{19.33}, tier=word]] [{NP} [{UT}, tier=word]]]]]]] [{VP} [{VP} [{suggests}, tier=word]] [{S} [{SCONJ} [{that}, tier=word]] [{S} [{NP} [{this}, tier=word]] [{VP} [{AUX} [{is}, tier=word]] [{NP} [{NP} [{DET} [{a}, tier=word]] [{NOUN} [{type-Ia}, tier=word]]] [{PP} [{ADP} [{at}, tier=word]] [{NP} [{NOUN} [{redshift}, tier=word]] [{NP} [{NOUN} [{z}, tier=word]] [{NP} [{0.044}, tier=word]]]]]]]]]]] [{PUNCT} [{.}, tier=word]]] \end{forest}\end{adjustbox}\\\small S6
\caption{Expected trees for the scientific sentences.}\label{fig:scientific}
\end{figure}

%==========================================================================
\section{Automatic analysis}

\subsection{PoS tagging with spaCy and sciSpaCy}
We used spaCy 3.7.5 with the model \texttt{en\_core\_web\_sm} 3.7.1 (trained on general English), and sciSpaCy with the model \texttt{en\_core\_sci\_sm} 0.5.4 (trained on biomedical texts). Table~\ref{tab:acc} compares them with our manual tags, and Table~\ref{tab:diff} lists every token where a model disagrees with us. When a model splits one of our tokens, we count it as an error.

\begin{table}[h]\centering\small
\caption{Accuracy against our manual tags.}\label{tab:acc}
\begin{tabular}{lccc}\toprule
Model & General (S1--S4) & Scientific (S5--S6) & All \\\midrule
\tinput{gen/pos_acc.tex}
\bottomrule\end{tabular}
\end{table}

\begin{table}[h]\centering\small
\caption{Tokens where a model disagrees with our manual tags.}\label{tab:diff}
\begin{tabularx}{\linewidth}{lll>{\raggedright\arraybackslash}X>{\raggedright\arraybackslash}X}\toprule
 & Token & Manual & \texttt{en\_core\_web\_sm} & \texttt{en\_core\_sci\_sm} \\\midrule
\tinput{gen/pos_diff.tex}
\bottomrule\end{tabularx}
\end{table}

\noindent Our main remarks:
\begin{itemize}
  \item On the general sentences both models are almost perfect. sciSpaCy tags \emph{flies} as NOUN, which is a possible reading, but with \emph{like} tagged ADP the sentence then has no verb.
  \item On the scientific sentences both models make the same mistakes on \emph{longer} (ADV) and \emph{hypocotyls} (ADJ): \emph{hypocotyls} is a rare plant-biology word. sciSpaCy is not better here, probably because it was trained on biomedical texts, not on plant biology or astronomy.
  \item Both tokenisers split \emph{J10354824+3900279}, and spaCy also splits \emph{type-Ia}.
  \item Some disagreements are conventions rather than errors. Both models tag \emph{Arabidopsis thaliana} as NOUN. For \emph{which} and \emph{this}, both models predict the same Penn Treebank tags (WDT, DT), but the conversion to universal tags gives PRON in spaCy and DET in sciSpaCy.
\end{itemize}

\subsection{Constituency parsing with CYK}
The provided \texttt{cyk.py} takes a list of possible tags for each word and returns all the trees that cover the sentence. It only accepts rules of the form X $\to$ Y Z and X $\to$ tag (Chomsky normal form), so one-word phrases are written as lexical rules: a NOUN, PROPN, PRON or NUM can directly be an NP, and a VERB can directly be a VP. We only keep the trees whose root is S. We wrote two grammars:
\begin{itemize}
  \item \textbf{G1} (general sentences): S $\to$ NP VP, S $\to$ S PUNCT, NP $\to$ DET NOUN, NP $\to$ NOUN NP, NP $\to$ NP PP, PP $\to$ ADP NP, VP $\to$ VERB NP, VP $\to$ VP PP\unskip.
  \item \textbf{G2} = G1 + the rules needed for the scientific sentences: NP $\to$ ADJ NP, NP $\to$ PROPN NP, NP $\to$ NUM NP, NP $\to$ NP VP, VP $\to$ AUX VERB, VP $\to$ AUX NP, VP $\to$ VP S, S $\to$ SCONJ S, S $\to$ PUNCT S, S $\to$ S S\unskip.
\end{itemize}
We ran CYK with four inputs: our manual tags, the spaCy tags, the sciSpaCy tags, and ``multiple PoS'', where each word receives our tag, its other possible PoS (Section~1.2) and the tags given by the two models. Table~\ref{tab:cyk} gives the number of trees.

\begin{table}[h]\centering\small
\caption{Number of trees returned by CYK (\checkmark: our expected tree is one of them).}\label{tab:cyk}
\begin{tabular}{llcccc}\toprule
 & Grammar & Manual tags & spaCy tags & sciSpaCy tags & Multiple PoS \\\midrule
\tinput{gen/cyk_counts.tex}
\bottomrule\end{tabular}
\end{table}

\paragraph{Time flies like an arrow.} With multiple PoS, G1 returns two trees: the expected one (Fig.~\ref{fig:timeflies}a) and ``time-flies like an arrow'' (Fig.~\ref{fig:timeflies}c). If we add two rules for imperative sentences (S $\to$ VERB NP and S $\to$ VP PP), we also get the two readings where \emph{Time} is a verb (b and d). With G2, the rule NP $\to$ NP VP adds a third tree meaning ``time [that] flies like an arrow'', which is wrong: the grammar does not check agreement. With the sciSpaCy tags there is no tree at all, because the sentence has no verb.

\begin{figure}[h]\centering
\begin{minipage}[b]{0.49\linewidth}\centering \begin{adjustbox}{max width=\linewidth}\begin{forest} for tree={s sep=1.2mm, l sep=2pt, l=0, inner sep=1pt, font=\footnotesize}, where n children=0{font=\footnotesize\itshape}{} [{S} [{NP} [{Time}, tier=word]] [{VP} [{VP} [{flies}, tier=word]] [{PP} [{ADP} [{like}, tier=word]] [{NP} [{DET} [{an}, tier=word]] [{NOUN} [{arrow}, tier=word]]]]]] \end{forest}\end{adjustbox}\\[-2pt]\small (a)\end{minipage}\hfill
\begin{minipage}[b]{0.49\linewidth}\centering \begin{adjustbox}{max width=\linewidth}\begin{forest} for tree={s sep=1.2mm, l sep=2pt, l=0, inner sep=1pt, font=\footnotesize}, where n children=0{font=\footnotesize\itshape}{} [{S} [{VERB} [{Time}, tier=word]] [{NP} [{NP} [{flies}, tier=word]] [{PP} [{ADP} [{like}, tier=word]] [{NP} [{DET} [{an}, tier=word]] [{NOUN} [{arrow}, tier=word]]]]]] \end{forest}\end{adjustbox}\\[-2pt]\small (b)\end{minipage}\\[6pt]
\begin{minipage}[b]{0.49\linewidth}\centering \begin{adjustbox}{max width=\linewidth}\begin{forest} for tree={s sep=1.2mm, l sep=2pt, l=0, inner sep=1pt, font=\footnotesize}, where n children=0{font=\footnotesize\itshape}{} [{S} [{NP} [{NOUN} [{Time}, tier=word]] [{NP} [{flies}, tier=word]]] [{VP} [{VERB} [{like}, tier=word]] [{NP} [{DET} [{an}, tier=word]] [{NOUN} [{arrow}, tier=word]]]]] \end{forest}\end{adjustbox}\\[-2pt]\small (c)\end{minipage}\hfill
\begin{minipage}[b]{0.49\linewidth}\centering \begin{adjustbox}{max width=\linewidth}\begin{forest} for tree={s sep=1.2mm, l sep=2pt, l=0, inner sep=1pt, font=\footnotesize}, where n children=0{font=\footnotesize\itshape}{} [{S} [{VP} [{VERB} [{Time}, tier=word]] [{NP} [{flies}, tier=word]]] [{PP} [{ADP} [{like}, tier=word]] [{NP} [{DET} [{an}, tier=word]] [{NOUN} [{arrow}, tier=word]]]]] \end{forest}\end{adjustbox}\\[-2pt]\small (d)\end{minipage}\\[6pt]
\caption{The four trees returned for \emph{Time flies like an arrow} with multiple PoS and the imperative rules.}\label{fig:timeflies}
\end{figure}

\paragraph{The spy sentences.} For S3 and S4, CYK always returns the same two trees, one with the PP attached to the verb phrase and one with the PP attached to the noun phrase (Fig.~\ref{fig:general}). The grammar only sees PoS tags, and the two sentences have the same tags, so it cannot choose the verb reading for the telescope and the noun reading for the revolver. This would require information about the words themselves, for example probabilities learned from a treebank.

\paragraph{The scientific sentences.} G1 cannot parse S5 and S6. With G2 our expected trees are found, but among 97 trees for S5 and 300 for S6: the grammar cannot decide where the \emph{when} clause, the relative clause, \emph{obtained on\ldots} or the PPs are attached, and the comma and the final period can also be attached at different levels. With the tags of spaCy or sciSpaCy there is no tree, because some tags (\emph{longer}/ADV, \emph{hypocotyls}/ADJ, \emph{this}/DET) and the extra tokens created by the tokenisers (\emph{+}, \emph{-}) cannot be combined by the grammar.

\subsection{Comparison with the Berkeley Neural Parser}
We used the Berkeley Neural Parser (benepar 0.2.0, model \texttt{benepar\_en3}) through its spaCy plugin. Compared with our expected trees:
\begin{itemize}
  \item S1, S2 and S3 have the same structure (S2 with the ``time passes'' reading, S3 with the PP attached to the verb).
  \item In S4, \emph{with the revolver} is also attached to the verb, as in S3.
  \item In S5, the \emph{when} clause is attached to the verb phrase as in our tree, but the relative clause is attached to \emph{high ambient temperature}. The parser's own tagger gets \emph{longer} and \emph{hypocotyls} right.
  \item In S6, the main structure is the same as ours, except that \emph{at redshift z 0.044} is attached to the verb phrase. The parser gives the same tree with spaCy's tokenisation, where CYK gave no tree.
\end{itemize}

%==========================================================================
\section{Observations}

\paragraph{PoS tagging.} Manual annotation is the most reliable because we can use the meaning of the sentence, but it is slow and depends on conventions, and the imposed tag set has no PRON. spaCy is fast and very accurate on general English, but it makes mistakes on unknown scientific words and splits identifiers. sciSpaCy did not do better on our scientific sentences, because its training domain (biomedical texts) is different from plant biology and astronomy. With only 73 tokens, these results are only indicative.

\paragraph{CYK with a context-free grammar.} Writing the grammar gives full control over the trees, but we saw several limits:
\begin{itemize}
  \item A rule added for one sentence can create new, wrong trees for other sentences (S2 with G2).
  \item CYK returns all possible trees and cannot rank them: S5 and S6 have hundreds of trees. A probabilistic grammar would be needed to choose the best one.
  \item The grammar only sees PoS tags, so it cannot use the meaning of words (telescope vs.\ revolver) or check agreement.
  \item One wrong tag or token is enough to get no tree at all.
  \item With only S, NP, VP and PP, subordinate clauses must be labelled S, and compounds such as \emph{Arabidopsis thaliana seedlings} can only be built right-branching.
\end{itemize}

\paragraph{Neural parser.} The Berkeley parser always returns a tree, handles unknown words and different tokenisations, and agrees with most of our choices. However, it returns a single tree even when the sentence is ambiguous, and its attachment choices (S4, the relative clause in S5) seem to follow what is frequent in its training data rather than the meaning of the sentence.

\end{document}
